% ============================================================
\section{Literature Review}
\label{sec:literature-review}
% ============================================================

This section surveys the bodies of work that inform the design of the
\textit{Ad LLM as a Service} platform. The review is organised by
\emph{approach} rather than by chronology: each subsection isolates a
distinct technical or methodological strand---digital advertising
formats, large language model capabilities, retrieval-augmented
generation, conversational-interface monetisation, prompt engineering
for controlled output, and engagement measurement---and positions the
present project relative to the state of the art in that strand.
Section~\ref{subsec:lit-summary} synthesises the findings and
identifies the specific research gap that this project addresses.

% ------------------------------------------------------------
\subsection{Evolution of Digital Advertising}
\label{subsec:evolution-ads}
% ------------------------------------------------------------

The economic architecture of online advertising has undergone three
major structural transitions since the first banner advertisement
appeared on HotWired in 1994, an AT\&T creative unit that reportedly
achieved a click-through rate exceeding 40\%~\cite{hotwired1994}.

\paragraph{From display to programmatic.}
The first generation of digital advertising was transactional and
manual: publishers sold rectangular display regions at fixed rates per
thousand impressions, and advertisers negotiated placements directly
with individual sites. The introduction of ad servers---most notably
DoubleClick in 1996---centralised delivery and measurement, but the
buying process remained fundamentally a media-planning exercise. The
emergence of real-time bidding (RTB) between 2007 and 2010 transformed
this model into an auction-based marketplace in which individual
impressions are priced algorithmically on the basis of predicted user
value~\cite{yuan2013rtb}. Demand-side platforms (DSPs), supply-side
platforms (SSPs), and data management platforms (DMPs) together created
a programmatic stack that, by 2022, accounted for approximately 84\%
of global digital display spending~\cite{emarketer2022programmatic}.

\paragraph{Search advertising and intent-based targeting.}
Google's introduction of AdWords in 2000, later renamed Google Ads,
established a second paradigm: intent-based targeting. Rather than
inferring user interest from page context or behavioural profiles,
search advertising matches sponsored results directly to the user's
expressed query. This alignment of commercial supply with declared
demand yields conversion rates substantially higher than those of
display advertising and has made search the single largest category of
digital ad spending. In 2024, Google's total advertising revenue
reached approximately \$265~billion, with the majority derived from
search placements~\cite{statista2024google}. The economic logic is
straightforward: a user who searches for ``best running shoes for flat
feet'' has explicitly declared a purchase-adjacent intent that no
amount of behavioural inference on a display network can reliably
reconstruct.

\paragraph{Social and behavioural targeting.}
A third strand developed in parallel on social-media platforms, where
dense self-reported demographic data and interaction graphs enabled
highly granular audience segmentation. Facebook (now Meta), Instagram,
and TikTok built advertising businesses on the principle that a
user's social graph, expressed interests, and engagement patterns
constitute a rich behavioural profile against which ads can be
targeted even in the absence of an explicit query. This model trades
the precision of declared intent for the breadth of passive
behavioural signal and has proven effective for awareness and
consideration campaigns, though it depends on third-party cookies and
device identifiers that are now being phased out under privacy
regulations such as the General Data Protection Regulation (GDPR) and
Apple's App Tracking Transparency framework~\cite{gdpr2016,
appleatt2021}.

\paragraph{The move toward native formats.}
Across all three paradigms, a persistent trend has been the migration
from visually distinct ad units---banners, interstitials,
pop-ups---toward formats that are structurally integrated with the
surrounding content. Native advertising, defined as paid content that
matches the form and function of the host medium, arose as a response
to banner blindness and ad-blocker adoption~\cite{wojdynski2016native}.
Sponsored articles on news sites, promoted posts on social feeds, and
influencer endorsements on video platforms all share the property that
the advertisement is embedded within the content stream rather than
placed adjacent to it. The platform examined in this report extends
this principle to its logical terminus in conversational AI: the
sponsored recommendation is not merely placed alongside the response
but woven into it by the language model itself.

% ------------------------------------------------------------
\subsection{Large Language Models and Conversational AI}
\label{subsec:llm-background}
% ------------------------------------------------------------

\paragraph{The transformer architecture.}
The architectural foundation of the present system is the transformer,
introduced by Vaswani et al.~\cite{vaswani2017attention}. The
transformer replaces the recurrent processing of earlier
sequence-to-sequence models with a self-attention mechanism that
computes pairwise relevance scores across all positions in the input
simultaneously, enabling parallelised training and effective capture
of long-range dependencies. This architecture has become the dominant
substrate for language modelling: every major large language model
(LLM) deployed as of 2025---including OpenAI's GPT
series~\cite{brown2020gpt3, openai2023gpt4}, Anthropic's Claude
family~\cite{anthropic2024claude3}, Google's Gemini~\cite{google2024gemini},
and Meta's LLaMA~\cite{touvron2023llama}---is built on transformer
variants. The progression from GPT-2's 117 million parameters to
GPT-4's estimated hundreds of billions illustrates the scaling
trajectory that has yielded the emergent capabilities exploited in
this project~\cite{wei2022emergent}.

\paragraph{Capabilities relevant to this project.}
Four capabilities of modern LLMs are directly exercised by the Ad LLM
platform:

\begin{enumerate}[label=(\roman*), leftmargin=2.5em]
  \item \emph{Instruction following.} Post-training alignment
    techniques---reinforcement learning from human feedback
    (RLHF)~\cite{ouyang2022instructgpt} and constitutional
    AI~\cite{bai2022constitutional}---have produced models that
    reliably follow detailed natural-language instructions, including
    multi-constraint prompts that specify output format, tone, and
    content restrictions.

  \item \emph{Structured output generation.} When instructed
    appropriately, LLMs generate syntactically valid JSON, markdown,
    or other structured formats with high consistency. This capability
    underpins the context extraction module, which must return a
    machine-parseable classification of intent, topics, sentiment, and
    ad receptivity from each user message.

  \item \emph{Contextual reasoning over long inputs.} Context windows
    of 100\,000 tokens and above---available in models such as
    Claude~3 (200\,000 tokens) and Gemini~1.5 (up to 1\,000\,000
    tokens)~\cite{anthropic2024claude3, google2024gemini}---allow the
    model to condition its generation on extensive conversation
    histories, a prerequisite for maintaining coherent multi-turn
    advertising integration.

  \item \emph{Tone adaptation.} LLMs modulate register, formality, and
    lexical choice in response to system-prompt instructions, enabling
    the response synthesis module to match the conversational style of
    the user while integrating a sponsored product mention.
\end{enumerate}

\paragraph{Lightweight versus capable models.}
A practical distinction that shapes the system architecture is the
trade-off between model capability and inference cost. Lightweight
models---exemplified by Anthropic's Haiku tier---offer low latency
and low per-token cost at the expense of reduced reasoning depth.
Capable models---exemplified by the Sonnet tier---produce higher-quality,
more nuanced text but at greater cost and latency. The Ad LLM
platform exploits this spectrum deliberately: context extraction, a
classification task requiring structured output but not creative
generation, is assigned to a Haiku-class model; response synthesis,
which demands natural integration of a product mention into a helpful
answer, is assigned to a Sonnet-class
model~\cite{anthropic2024claude3}. This two-tier strategy is
consistent with the broader industry practice of routing sub-tasks to
the smallest model that achieves acceptable quality, thereby
minimising aggregate inference cost~\cite{chen2023frugalgpt}.

% ------------------------------------------------------------
\subsection{Retrieval-Augmented Generation (RAG)}
\label{subsec:rag}
% ------------------------------------------------------------

\paragraph{The RAG paradigm.}
Retrieval-augmented generation, formalised by Lewis et
al.~\cite{lewis2020rag}, addresses a fundamental limitation of
parametric language models: their knowledge is frozen at training time
and cannot be updated or extended without retraining. The RAG
architecture augments the generative model with a non-parametric
memory---a corpus of documents indexed for efficient retrieval---so
that, at inference time, the model conditions its output on both the
user query and the most relevant retrieved passages. Lewis et al.\
demonstrated that this combination of a pre-trained seq2seq
transformer (BART) with a dense vector index of Wikipedia, accessed
via a neural retriever (Dense Passage Retriever,
DPR)~\cite{karpukhin2020dpr}, achieved state-of-the-art results on
open-domain question answering benchmarks and generated more factual,
specific language than parametric-only baselines.

\paragraph{Vector embeddings and semantic similarity.}
The retrieval component of a RAG pipeline depends on dense vector
representations of both the query and the candidate documents.
Embedding models---such as OpenAI's \texttt{text-embedding-3-small},
Cohere's Embed, and open-source alternatives like BGE and
E5~\cite{wang2022e5, xiao2023bge}---map text into high-dimensional
vector spaces in which semantic similarity corresponds to geometric
proximity (typically measured by cosine distance). The quality of
these embeddings determines retrieval precision: a query embedding
that accurately captures the user's informational need will surface
documents whose content is genuinely relevant, while a poor embedding
will introduce noise that degrades the generator's output.

\paragraph{Vector databases and extensions.}
At scale, approximate nearest-neighbour (ANN) search over embedding
vectors requires specialised indexing. Dedicated vector databases such
as Pinecone, Weaviate, and Milvus provide managed infrastructure for
this purpose~\cite{jing2024vectordb}. For applications with moderate
inventory sizes---tens of thousands of records rather than
billions---the \texttt{pgvector} extension to PostgreSQL offers a
pragmatic alternative: it adds vector-typed columns and
cosine-distance operators to a relational database, avoiding the
operational overhead of a separate vector
store~\cite{pgvector2023}. The Ad LLM platform adopts
\texttt{pgvector} for precisely this reason: the ad inventory is
expected to contain at most tens of thousands of items, and
co-locating vector search with the relational schema for budget
management, frequency capping, and event logging simplifies both
development and deployment.

\paragraph{Ad matching as a specialised RAG pipeline.}
The conceptual contribution of the present project's ad-matching
subsystem is the reinterpretation of RAG in a commercial context. In
a conventional RAG system, the retrieval step fetches passages that
are informationally relevant to the user's question; in the Ad LLM
pipeline, the retrieval step fetches \emph{advertisements} that are
semantically relevant to the user's conversational context. The query
embedding is constructed not from the raw user message but from the
structured context extracted by the lightweight LLM---topics, intent,
and entities---ensuring that the retrieval signal reflects interpreted
user need rather than surface lexical overlap. A business-rule
re-ranking stage then applies non-semantic constraints (budget
availability, bid price, per-session frequency caps) before the top
candidates are passed to the generator. This architecture preserves
the retrieve-then-generate structure of RAG while specialising both
the retrieval objective and the post-retrieval filtering to the
requirements of advertising.

% ------------------------------------------------------------
\subsection{Advertising in Conversational Interfaces}
\label{subsec:conversational-ads}
% ------------------------------------------------------------

\paragraph{Industry deployments.}
The integration of advertising into conversational AI is a recent and
rapidly evolving development. Microsoft was the first major platform
to insert sponsored content into a conversational interface, embedding
text ads, shopping campaigns, and multimedia placements into Bing Chat
responses shortly after its launch in February
2023~\cite{microsoft2023bingads}. In September 2023, Microsoft
announced ``Conversational Ads,'' a new format designed specifically
for the chat medium, including ``Compare \& Decide Ads'' that present
structured product comparisons within the dialogue
flow~\cite{microsoft2023conversationalads}. Google followed with
experiments in its Search Generative Experience (SGE), placing
sponsored results both adjacent to and within AI-generated answer
summaries, labelled with a bold ``Sponsored''
tag~\cite{google2023sge}. By 2025, Google had extended these
placements into AI Mode, a fully conversational search interface
powered by Gemini~2.0, and reported that ads drawn from existing
Search and Shopping campaigns were being served within AI-generated
responses~\cite{google2025aimode}.

These deployments share a common limitation from the perspective of
the present project: they append or interleave conventional ad units
(text links, product cards) alongside AI-generated text rather than
instructing the language model to integrate the product mention
\emph{into} the response itself. The result is a hybrid layout in
which the advertisement remains visually and linguistically distinct
from the organic answer---a format that is arguably native to the
search-results page but not to the conversational medium.

\paragraph{Academic work on native advertising and user perception.}
A substantial body of research examines how users perceive and respond
to native advertising. Wojdynski and Evans~\cite{wojdynski2016native}
demonstrated that fewer than one in ten participants recognised
article-style native advertising as paid content, even when disclosure
labels were present, highlighting the tension between format
integration and consumer awareness. Subsequent work has shown that
this tension operates through two competing
mechanisms~\cite{vandam2022disclosure}: explicit sponsorship
disclosures activate persuasion knowledge, prompting critical
evaluation and reducing brand attitude, but simultaneously enhance
perceived transparency, which increases trust and partially offsets
the negative effect. The balance between these mechanisms depends on
disclosure language, placement, and
prominence~\cite{wojdynski2016native, campbell2018disclosure}. A
Nielsen study reported that 79\% of users correctly identified content
as advertising when the label read ``Advertisement,'' compared to only
48\% for the label ``Presented
by''~\cite{nielsen2014nativead}. These findings directly inform the
design decision in the Ad LLM platform to use a mandatory, visually
rendered \texttt{[Sponsored]} tag---a disclosure format that is both
explicit and machine-enforceable.

\paragraph{Ethical considerations: transparency and trust.}
The ethical literature on advertising transparency converges on a
central finding: consumers respond more favourably to advertising when
they perceive that the advertiser is being honest about the commercial
nature of the content~\cite{carr2014disclosure, campbell2018disclosure}.
Transparency functions as a signal of brand
integrity~\cite{kang2019transparency}, and the downstream effects on
trust and purchase intention are robust across experimental
replications. In the conversational AI context, the stakes are
amplified: the user has entered a dialogue that they expect to be
helpful and impartial, and undisclosed commercial influence would
constitute a violation of that implicit contract. The Ad LLM platform
addresses this risk through three mechanisms: (i)~the
\texttt{[Sponsored]} tag is mandatory and enforced at the prompt
level; (ii)~the response-generation prompt instructs the model to
ignore the advertisement if it is not relevant to the user's question;
and (iii)~conversations classified as sensitive are automatically
excluded from ad injection.

\paragraph{Regulatory frameworks.}
The Federal Trade Commission (FTC) revised its Endorsement Guides in
July 2023, updating the definition of ``endorsement'' to encompass
digital formats including social-media tags, virtual influencers, and
AI-generated recommendations~\cite{ftc2023endorsement}. The revised
Guides require that any material connection between an advertiser and
an endorser be disclosed ``clearly and conspicuously,'' using language
that is difficult to miss and easily understood by the target
audience. The Guides explicitly note that platform-provided disclosure
tools (e.g., Instagram's ``Paid partnership'' label) may not be
sufficient on their own. Although the FTC has not yet issued
guidance specific to LLM-generated sponsored content, the general
principle---that consumers must be informed when content is
commercially motivated---applies directly. The \texttt{[Sponsored]}
disclosure mechanism in the Ad LLM platform is designed to satisfy
this principle by ensuring that every sponsored product mention is
preceded by an unambiguous, machine-enforced marker.

% ------------------------------------------------------------
\subsection{Prompt Engineering for Controlled Generation}
\label{subsec:prompt-engineering}
% ------------------------------------------------------------

\paragraph{Techniques for reliable structured output.}
Prompt engineering---the practice of crafting input instructions to
elicit desired model behaviour without modifying model
parameters---has emerged as a critical discipline in applied LLM
systems. Schulhoff et al.~\cite{schulhoff2024prompt} present a
taxonomy of 58~distinct prompting techniques, ranging from zero-shot
and few-shot prompting to chain-of-thought reasoning, self-refinement,
and retrieval-augmented prompting. Sahoo et
al.~\cite{sahoo2024prompt} provide a complementary survey organised
by application area, demonstrating that prompt design significantly
affects output quality across tasks including question answering, code
generation, and text classification. For the Ad LLM platform, the
most relevant techniques are \emph{role assignment} (instructing the
model to adopt a specific persona), \emph{output-format
specification} (requiring JSON or markdown structure), and
\emph{conditional instruction} (directing the model to include or
exclude content based on stated criteria).

\paragraph{System prompts and instruction following.}
The system prompt---a privileged instruction block that precedes the
user's message and persists across turns---is the primary mechanism
through which the Ad LLM platform controls the response-generation
model. The system prompt specifies: the model's role (a helpful
assistant), the conditions under which a sponsored product should be
mentioned (relevance to the user's question), the required disclosure
format (\texttt{[Sponsored]} tag), and the link format (markdown with
tracking URL). Ouyang et al.~\cite{ouyang2022instructgpt} showed that
instruction-tuned models follow system-level directives with high
fidelity, though compliance is not absolute: adversarial or ambiguous
inputs can cause the model to deviate from instructions. The prompt
design for the Ad LLM platform therefore includes explicit rejection
logic (``If the product is not relevant to the conversation, ignore it
entirely'') and is validated against an evaluation set of 20--30
scenarios covering both successful integration and correct rejection.

\paragraph{Balancing helpfulness with ad integration naturalness.}
A distinctive prompt-engineering challenge in this project is the
dual-objective nature of the generation task: the response must
simultaneously be \emph{helpful} (answering the user's question
accurately and completely) and \emph{commercially effective} (weaving
a product mention into the answer in a way that feels like an informed
suggestion rather than an insertion). These objectives are not
inherently conflicting---a relevant product recommendation \emph{is}
helpful---but they become so when the matched advertisement is only
tangentially related to the user's query. The prompt must therefore
encode a threshold: the model should integrate the product only when
doing so genuinely enhances the response, and should produce a
purely organic answer otherwise. Achieving consistent behaviour
across diverse conversation types requires iterative prompt refinement
guided by systematic evaluation.

\paragraph{LLM-as-judge for automated quality evaluation.}
Manual evaluation of generated responses does not scale. The
LLM-as-judge paradigm, surveyed by Zheng et
al.~\cite{zheng2023llmjudge} and Gu et al.~\cite{gu2024llmjudge},
addresses this limitation by using a separate LLM call to score
generated outputs along specified quality dimensions. Zheng et al.\
demonstrated that GPT-4, when used as a judge, achieves approximately
80\% agreement with human preference scores---matching the
inter-annotator agreement rate among human raters themselves. The Ad
LLM platform adopts this paradigm for offline quality monitoring: a
judge model scores sampled production responses on helpfulness,
naturalness of ad integration, disclosure accuracy, and ad relevance,
enabling automated detection of quality regressions without requiring
continuous human review.

% ------------------------------------------------------------
\subsection{User Engagement Metrics in Conversational Systems}
\label{subsec:engagement-metrics}
% ------------------------------------------------------------

\paragraph{Traditional ad metrics.}
The standard measurement framework for digital advertising centres on
three quantities: impressions (the number of times an ad is rendered),
click-through rate (CTR, the fraction of impressions that result in a
click), and conversion rate (the fraction of clicks that result in a
desired action, such as a purchase). These metrics are well-understood
and form the basis of pricing models including cost per mille (CPM)
and cost per click (CPC). However, they were designed for a visual
medium in which the user's attention is inferred from the rendering of
a display unit. In a conversational interface, where the entire
response is a continuous text stream, the concept of a discrete
``impression'' is less well defined, and a ``click'' on an embedded
link may not capture the full spectrum of user engagement with the
recommended product.

\paragraph{Novel conversational metrics.}
The Ad LLM platform introduces two metrics that exploit the unique
properties of the conversational medium:

\begin{enumerate}[label=(\roman*), leftmargin=2.5em]
  \item \emph{Follow-up engagement.} After an advertisement is
    included in a response, the system monitors the user's next
    message for references to the advertised product (e.g., ``Tell me
    more about the Pegasus~41''). A follow-up that engages with the
    product constitutes a stronger signal of genuine interest than a
    passive impression or even a click, because it requires the user
    to actively formulate a query about the product within the
    conversational context.

  \item \emph{Attention verification.} The time elapsed between
    response delivery and the user's next action provides a proxy for
    read-through. A response that is read carefully before the user
    continues the conversation suggests higher attentional engagement
    than one that is followed immediately by an unrelated query. This
    metric, while noisy, enables the construction of ``verified
    attention'' scores that weight impressions by estimated engagement
    depth---a concept analogous to viewability measurement in display
    advertising~\cite{ias2023viewability}.
\end{enumerate}

\paragraph{Why conversational engagement is a stronger signal.}
The fundamental advantage of conversational engagement metrics over
traditional click-through measurement is that they capture
\emph{active cognitive processing} rather than a binary
click-or-not-click signal. A user who asks a follow-up question about
an advertised product has read the recommendation, understood it,
found it relevant, and invested the effort to formulate a new query.
This chain of cognitive steps is substantially more informative than a
click, which may be accidental, curiosity-driven, or immediately
followed by a bounce. Conversational engagement therefore offers
advertisers a higher-fidelity signal of genuine interest, which in
turn supports more accurate attribution and, ultimately, more
efficient allocation of advertising spend.

\paragraph{Related work on measuring user attention in dialogue.}
Research on dialogue systems has developed several proxies for user
engagement, including turn length, response latency, sentiment
trajectory, and topic persistence~\cite{see2019engagement,
mehri2020unsupervised}. See et al.~\cite{see2019engagement} found
that longer user responses and sustained topic continuity correlate
with self-reported satisfaction in open-domain dialogue. Mehri and
Esk{\'e}nazi~\cite{mehri2020unsupervised} proposed unsupervised
dialogue quality metrics based on response coherence and relevance
scoring. The engagement classifier in the Ad LLM platform draws on
this work by treating topic persistence (continued discussion of the
advertised product) as the primary engagement signal, initially
implemented via keyword and entity matching with a planned upgrade to
a lightweight LLM-based classifier as labelled data becomes available.

% ------------------------------------------------------------
\subsection{Summary and Research Gap}
\label{subsec:lit-summary}
% ------------------------------------------------------------

The preceding subsections have surveyed five bodies of work that
converge on the problem addressed by this project: (i)~the evolution
of digital advertising toward native, content-integrated formats;
(ii)~the capabilities of transformer-based LLMs for instruction
following, structured output, and tone-adapted generation; (iii)~the
RAG paradigm for grounding language-model output in retrieved external
context; (iv)~the emerging practice of embedding advertising in
conversational AI interfaces, together with the ethical and regulatory
constraints that govern disclosure; (v)~prompt engineering techniques
for controlled generation and the LLM-as-judge paradigm for automated
quality evaluation; and (vi)~engagement metrics that go beyond
click-through to capture genuine user attention in dialogue systems.

The gap that this project addresses lies at the intersection of these
strands. Existing industry deployments of conversational
advertising---Microsoft's Bing Chat ads and Google's SGE/AI Mode
placements---insert conventional ad units (text links, product cards)
alongside AI-generated responses rather than instructing the language
model to integrate product mentions into the response itself.
Academic work on native advertising has extensively studied user
perception and disclosure effectiveness, but has not examined
LLM-generated native advertising in which the model itself produces
the sponsored content. The RAG literature provides the retrieval
infrastructure for context-grounded generation, but has not applied
the retrieve-then-generate pipeline to ad matching and integration.
Prompt engineering research offers the techniques for controlled
output, but has not specifically addressed the dual-objective problem
of simultaneous helpfulness and natural ad integration.

No open, well-documented architecture currently exists that combines
semantic ad matching via a specialised RAG pipeline, LLM-driven
natural ad integration with mandatory disclosure, sensitivity-aware ad
gating, and engagement-based measurement within a single end-to-end
system. The \textit{Ad LLM as a Service} platform is designed to
address this gap: it provides a complete, working pipeline from user
message to ad-augmented response, accompanied by an analytics
framework that measures conversational engagement as a first-class
advertising metric.